\documentclass[11pt]{article}
\usepackage{coursenotes}
\begin{document}
\handoutheader{H5}{Regression discontinuity}{Identification from a threshold rule}{Meeting 8 (Double machine learning)}

\begin{decision}
Many business rules switch treatment on at a threshold: a loyalty tier at 100 visits, free shipping above a basket value, a
credit limit at a score cut-off, a rating rounded up to four stars. Units just above and just below the threshold differ in
treatment but, if they cannot place themselves precisely, in little else. The jump in outcomes at the threshold is then an
effect, for units at the threshold.
\end{decision}

\begin{goals}
  \item state the sharp and fuzzy regression discontinuity (RD) designs and their identifying assumption (continuity);
  \item estimate the jump by local linear regression and choose a bandwidth;
  \item run the standard checks: density of the running variable, covariate balance, bandwidth, and placebo cut-offs;
  \item interpret the effect as local and say what it implies for moving the threshold.
\end{goals}

%=====================================================================================
\sectionone

\subsection{The sharp design}

Let $R$ be a \emph{running variable} and $c$ a cut-off, with $D = \ind{R \ge c}$. Treatment is a deterministic function of $R$, so
overlap fails everywhere: there is no unit at any $R$ with both treatments. Identification comes instead from comparing units on
either side of $c$.

\begin{assum}{Continuity}{cont}
$\E[Y(0) \mid R = r]$ and $\E[Y(1) \mid R = r]$ are continuous in $r$ at $c$.
\end{assum}

\begin{prop}{Sharp RD identification}{sharp}
Under consistency and continuity,
\[
  \tau_{\text{RD}} \equiv \lim_{r \downarrow c}\E[Y \mid R = r] - \lim_{r \uparrow c}\E[Y \mid R = r] = \E[Y(1) - Y(0) \mid R = c].
\]
\end{prop}
\begin{proof}
Above $c$, $\E[Y \mid R = r] = \E[Y(1) \mid R = r]$, whose limit as $r \downarrow c$ is $\E[Y(1) \mid R = c]$ by continuity; below $c$ the same
holds for $Y(0)$. Subtract.
\end{proof}

Continuity fails when units can sort precisely across the cut-off (members who log extra visits just to qualify) or when something else
also changes at $c$ (a different price list starts at the same basket value). The first can be partly checked in the data; the second must be
argued from how the business runs.

\subsection{The fuzzy design}

Often the threshold changes only the probability of treatment (the tier is offered at 100 visits, but members must opt in; some below get it by
promotion). Let $Z = \ind{R \ge c}$ act as an instrument at the cut-off.

The tool is the LATE theorem from Meeting~2. Take a randomized offer $Z$ and treatment received $D$, and call a unit a
\emph{complier} if it takes the treatment when offered and not otherwise. If the offer moves the outcome only through the
treatment (exclusion) and no unit takes the treatment only when not offered (monotonicity), then the ratio of the offer's effect on $Y$
(the intention-to-treat effect, ITT) to its effect on $D$ (the first stage) equals the local average treatment effect (LATE), the
average effect of treatment among compliers.

\begin{prop}{Fuzzy RD identification}{fuzzy}
If, near $c$, the conditions of the instrumental-variables LATE theorem hold with $Z = \ind{R \ge c}$ (continuity of potential outcomes and potential
treatments in $r$, monotonicity, and a nonzero jump in $\Pr(D = 1 \mid R)$), then
\[
  \frac{\lim_{r \downarrow c}\E[Y \mid r] - \lim_{r \uparrow c}\E[Y \mid r]}{\lim_{r \downarrow c}\E[D \mid r] - \lim_{r \uparrow c}\E[D \mid r]}
  = \E[Y(1) - Y(0) \mid \text{complier},\ R = c] .
\]
\end{prop}
\begin{proof}
At the cut-off, being just above or just below is as good as random by continuity, so the numerator is the ITT and the denominator the first stage
for units at $c$; apply the LATE theorem locally.
\end{proof}

The effect is local twice over: to compliers, and to units at the threshold.

\subsection{Estimation: local linear regression}

The limits are estimated by fitting a line on each side of $c$, using only units within a \emph{bandwidth} $h$ of the cut-off, usually weighting
units closer to $c$ more (a triangular kernel $w_i = \max\{0, 1 - |R_i - c|/h\}$):
\[
  Y_i = \alpha + \tau\,\ind{R_i \ge c} + \beta_0(R_i - c) + \beta_1(R_i - c)\ind{R_i \ge c} + \varepsilon_i, \quad |R_i - c| < h .
\]
$\hat\tau$ is the RD estimate. Lines rather than means within the window remove the bias from the slope of $\E[Y \mid R]$ near $c$. Avoid high-order global
polynomials, which put large weight on far-away units and can produce spurious jumps.

\paragraph{The bandwidth.} A small $h$ uses units closest to $c$ (less bias, more variance); a large $h$ the reverse. Data-driven bandwidths minimize
the estimated mean squared error (MSE). Because the MSE-optimal bandwidth leaves bias of the same order as the SE, conventional intervals undercover;
use robust bias-corrected intervals (Calonico et al. 2014), and show results across a range of bandwidths.

\subsection{Checks}

\begin{itemize}
  \item \textbf{Density of $R$ at $c$.} A pile-up of units just above the cut-off suggests manipulation. Compare counts in narrow bins on either side or
  use a formal density test.
  \item \textbf{Covariate balance.} Pre-determined covariates should not jump at $c$. Run the RD regression with each covariate as the outcome.
  \item \textbf{Bandwidth sensitivity.} The estimate should be stable across reasonable bandwidths.
  \item \textbf{Placebo cut-offs.} Estimate ``jumps'' at values of $R$ where nothing changes; they should be near zero.
\end{itemize}

\subsection{What a local effect supports}

$\tau_{\text{RD}}$ is the effect for units at the threshold. It directly supports small moves of the threshold (from 100 to 95 visits), if the effect is
roughly constant over the units the move would add. It says little about units far from $c$. A \emph{kink} design uses a change in the slope of a
rule (a subsidy that rises with spend up to a cap) in the same way, identifying a derivative rather than a level.

\begin{pitfall}[RD errors]
Fitting high-order polynomials over the whole range; reporting one bandwidth; ignoring bunching just above the cut-off; another rule changing at the
same threshold; reading the local effect as an average for everyone.
\end{pitfall}

%=====================================================================================

\sectiontwo

\paragraph{Setting.} FitLife grants ``Gold'' status, with perks costing ¥60 a year, to members with at least 100 visits in the year. Gold members renew
more often, but they are also the most committed. A renewal is worth ¥1{,}000, so Gold pays at the margin if it raises renewal by more than $0.06$.
Should the threshold be lowered to 95 visits? The data are 60{,}000 member-years from a simulation calibrated to this setting
(true jump $0.08$).

\paragraph{Step 1: naive.} Members at or above 100 visits renew $0.33$ more often than those below: mostly commitment, not Gold.

\begin{center}
\begin{tikzpicture}
\begin{axis}[width=0.7\linewidth, height=5cm, xlabel={visits in the year ($R$)}, ylabel={renewal rate}, axis lines=left,
  xmin=40, xmax=160, ymin=0.2, ymax=0.9]
\addplot[only marks, mark=*, mark size=1.2pt, noteblue] table[col sep=comma, x index=0, y index=1] {h5_bins.csv};
\draw[dashed, warnred] (axis cs:100,0.2) -- (axis cs:100,0.9);
\end{axis}
\end{tikzpicture}

\small Figure 1. Renewal by bins of five visits. The cut-off is dashed; the jump there is the effect of Gold for members at 100 visits.
\end{center}

\paragraph{Step 2: the RD estimate.} Local linear regression with a triangular kernel:
\begin{center}
\begin{tabular}{lccc}
\toprule
Bandwidth (visits) & Members used & $\hat\tau$ & SE \\
\midrule
10 & 9{,}375 & 0.091 & 0.022 \\
20 & 19{,}197 & 0.087 & 0.015 \\
30 & 28{,}983 & 0.084 & 0.012 \\
\bottomrule
\end{tabular}
\end{center}
The estimate is stable across bandwidths at about $0.085$, against a naive gap of $0.33$.

\paragraph{Step 3: checks.} Counts in the five visit-values below and above the cut-off are $2{,}503$ and $2{,}472$: no pile-up. Members' age does not
jump at the cut-off (estimate $0.30$ years, SE $0.25$).

\paragraph{Step 4: the decision.} At the threshold Gold raises renewal by about $0.085$; at $h = 20$ the 95\% interval is $[0.057, 0.117]$. The
point estimate clears the $0.06$ break-even, but the interval's lower end falls just below it. Lowering the threshold to 95 adds members just below the current cut-off; if their effect resembles that at 100, the move pays.
That is an assumption, mild for a five-visit move and not for a 50-visit one.

\begin{exercise}[computation]
Suppose Gold were opt-in: just above 100 visits, 65\% of members hold Gold; just below, 5\% do (via promotions). The jump in renewal at the cut-off is
$0.052$. What is the fuzzy-RD effect, and for whom?
\end{exercise}
\answer{$0.052/(0.65 - 0.05) = 0.087$: the effect of Gold on renewal for members at 100 visits who take Gold because they cross the threshold (compliers
at the cut-off).}

\begin{exercise}[judgment]
Marketing reveals that members get a reminder email (``you're 5 visits from Gold!'') at 95 visits. What does this do to the design?
\end{exercise}
\answer{Members who receive the reminder may add visits to qualify, so members just above 100 are more motivated than those just below: sorting
around the cut-off, which breaks continuity. Expect bunching just above 100; check the density. A cleaner comparison uses a running variable fixed
before the reminder (visits by a date before reminders start), or treats the reminder as part of the treatment.}

%=====================================================================================
\sectionthree

\subsection*{Annotated reading}
\begin{readinglist}
\reading{Cattaneo, Matias D., Nicol\'as Idrobo, and Roc\'io Titiunik. 2020. \emph{A Practical Introduction to Regression
Discontinuity Designs: Foundations}. Cambridge Elements: Quantitative and Computational Methods for the Social Sciences.
Cambridge: Cambridge University Press. \url{https://doi.org/10.1017/9781108684606}}{The modern practice: local polynomials,
bandwidths, robust inference, checks}{All}{2}
\reading{Imbens, Guido W., and Thomas Lemieux. 2008. ``Regression Discontinuity Designs: A Guide to Practice.''
\emph{Journal of Econometrics} 142 (2): 615--35.}{A concise practical guide}{All}{2}
\reading{Lee, David S., and Thomas Lemieux. 2010. ``Regression Discontinuity Designs in Economics.'' \emph{Journal of
Economic Literature} 48 (2): 281--355. \url{https://doi.org/10.1257/jel.48.2.281}}{Intuition, the as-good-as-random
interpretation, many examples}{Sections 1--4}{1}
\reading{Calonico, Sebastian, Matias D. Cattaneo, and Rocio Titiunik. 2014. ``Robust Nonparametric Confidence Intervals for
Regression-Discontinuity Designs.'' \emph{Econometrica} 82 (6): 2295--2326. \url{https://doi.org/10.3982/ECTA11757}}{Why
MSE-optimal bandwidths need bias correction}{Sections 1--3}{3}
\reading{Hahn, Jinyong, Petra Todd, and Wilbert van der Klaauw. 2001. ``Identification and Estimation of Treatment Effects
with a Regression-Discontinuity Design.'' \emph{Econometrica} 69 (1): 201--9.
\url{https://doi.org/10.1111/1468-0262.00183}}{The fuzzy design as local instrumental variables}
{Sections 1--3}{3}
\end{readinglist}

\end{document}
